\begin{lemma}
Let $H\colon\mathbb{C}\to\mathbb{C}$, let $s,k,t,l\ge0$ be integers, and let
$P\in\mathbb{P}(s,k)$ and $Q\in\mathbb{P}(t,l)$ relative to $H$ in the sense
of \noderef{PClass}.  Then the pointwise product $PQ$ belongs to
$\mathbb{P}(s+t,k+l)$ relative to $H$.
\end{lemma}
\begin{proof}
By \noderef{PClass} applied to $P$, choose an integer $r\ge0$, complex numbers
$c_1,\dots,c_r$ and lists $M_1,\dots,M_r$ of derivative symbols with
\begin{equation}\tag{1}
  s(M_\alpha)=s,\qquad k(M_\alpha)=k\qquad(\alpha=1,\dots,r),
\end{equation}
\begin{equation}\tag{2}
  P(\xi)=\sum_{\alpha=1}^{r}c_\alpha H^{M_\alpha}(\xi)\qquad\text{for every }\xi\in\mathbb{C},
\end{equation}
and by \noderef{PClass} applied to $Q$, choose an integer $q\ge0$, complex
numbers $b_1,\dots,b_q$ and lists $N_1,\dots,N_q$ with
\begin{equation}\tag{3}
  s(N_\beta)=t,\qquad k(N_\beta)=l\qquad(\beta=1,\dots,q),
\end{equation}
\begin{equation}\tag{4}
  Q(\xi)=\sum_{\beta=1}^{q}b_\beta H^{N_\beta}(\xi)\qquad\text{for every }\xi\in\mathbb{C}.
\end{equation}

Consider the $rq$ pairs $(\alpha,\beta)$ with $1\le\alpha\le r$,
$1\le\beta\le q$.  To each we attach the complex number $c_\alpha b_\beta$ and
the list
\[
  M_\alpha * N_\beta
\]
obtained by concatenating $M_\alpha$ and $N_\beta$, i.e.\ by listing the
entries of $M_\alpha$ followed by the entries of $N_\beta$.

\emph{Step 1: every attached list has type $(s+t,k+l)$.}  By
\noderef{MonomialValue}, $k(\,\cdot\,)$ is the length of a list and
$s(\,\cdot\,)$ is the sum of the orders of its entries.  Concatenation lists
the entries of the two lists one after the other, so its length is the sum of
the two lengths and the sum of the orders of its entries is the sum of the two
sums:
\[
  k(M_\alpha * N_\beta)=k(M_\alpha)+k(N_\beta)=k+l,
  \qquad
  s(M_\alpha * N_\beta)=s(M_\alpha)+s(N_\beta)=s+t,
\]
using (1) and (3).

\emph{Step 2: the monomial of a concatenation is the product of the
monomials.}  Again by \noderef{MonomialValue}, $H^{M}(\xi)$ is the product of
the symbol values $H^{d}(\xi)$ over the entries $d$ of $M$.  For a
concatenation the list of entries is the concatenation of the two lists of
entries, so the product splits as the product of the two partial products:
\begin{equation}\tag{5}
  H^{M_\alpha * N_\beta}(\xi)=H^{M_\alpha}(\xi)\,H^{N_\beta}(\xi)
  \qquad\text{for every }\xi\in\mathbb{C}.
\end{equation}

\emph{Step 3: the linear combination.}  Let $\xi\in\mathbb{C}$.  Using (5) and
then the distributive law for finite sums of complex numbers,
\[
  \sum_{\alpha=1}^{r}\sum_{\beta=1}^{q}c_\alpha b_\beta\,H^{M_\alpha * N_\beta}(\xi)
  =\sum_{\alpha=1}^{r}\sum_{\beta=1}^{q}
     \bigl(c_\alpha H^{M_\alpha}(\xi)\bigr)\bigl(b_\beta H^{N_\beta}(\xi)\bigr)
  =\Bigl(\sum_{\alpha=1}^{r}c_\alpha H^{M_\alpha}(\xi)\Bigr)
   \Bigl(\sum_{\beta=1}^{q}b_\beta H^{N_\beta}(\xi)\Bigr),
\]
and by (2) and (4) the right-hand side is $P(\xi)Q(\xi)$.

The finite family of coefficients $c_\alpha b_\beta$ together with the lists
$M_\alpha * N_\beta$, which all have type $(s+t,k+l)$ by Step 1, therefore
presents $PQ$ as a finite $\mathbb{C}$-linear combination of monomials of type
$(s+t,k+l)$, valid at every point of $\mathbb{C}$.  By \noderef{PClass} this
says exactly that $PQ\in\mathbb{P}(s+t,k+l)$ relative to $H$.
\end{proof}
